\labsection{2}{Heuristic Search on the Airline Network}{sec:lab02-informed}

\begin{labproject}{Compare UCS, Greedy Best-First Search, and A$^\star$}
\textbf{Scenario.} Improve the Lab 1 planner by adding geographic knowledge while preserving a clear quality baseline.\\
\textbf{Goal.} Implement GBFS and A$^\star$, validate the Haversine heuristic, and measure the trade-off between solution cost and search effort.

\begin{labmode}
Reuse the Lab 1 OpenFlights files and run \path{all_experiments/lab02_informed_search_worked.ipynb}. Save evidence under \texttt{results/lab02/}.
\end{labmode}

\begin{mlsteps}
  \item \textbf{Establish the baseline.} Run UCS for the three standard route queries and save its optimal route cost.
  \item \textbf{Implement the heuristic.} Use Haversine distance from the current airport to the destination. Implement GBFS with $h(n)$ and A$^\star$ with $g(n)+h(n)$.
  \item \textbf{Check admissibility empirically.} Compare $h(n)$ against exact route cost-to-go for at least 10 reachable airports per query.
  \item \textbf{Compare methods.} Record path cost, expanded nodes, maximum frontier, and runtime for UCS, GBFS, and A$^\star$.
  \item \textbf{Stress the guarantee.} Repeat the queries with an inflated heuristic $1.5h(n)$ and explain when the returned cost changes.
  \item \textbf{Save evidence.} Export the algorithm metrics, heuristic check, and an expanded-nodes figure.
\end{mlsteps}

\requiredoutput{Executed notebook, \texttt{metrics.csv}, \texttt{heuristic\_check.csv}, and search-effort figure.}
\end{labproject}

\begin{verification}
With the ordinary Haversine heuristic, A$^\star$ must match UCS route cost on the tested queries. Investigate priority updates and closed-set handling if it does not.
\end{verification}

\begin{labdeliverable}
Push the completed Python work and create tag \texttt{lab02-final}. The README must state whether A$^\star$ preserved UCS cost, how much expansion was reduced, and where GBFS returned a higher-cost route.
\end{labdeliverable}
